\documentclass[11pt,letterpaper]{article}
\usepackage[T1]{fontenc}
\usepackage[margin=0.87in,top=0.86in,bottom=0.76in,headheight=23pt,headsep=15pt,footskip=28pt]{geometry}
\usepackage{amsmath,amssymb}
\usepackage{newpxtext,newpxmath}
\usepackage[semibold]{sourcesans}
\usepackage[varqu]{inconsolata}
\usepackage{microtype,xcolor,enumitem,fancyhdr}
\usepackage[most]{tcolorbox}
\usepackage{hyperref}
\definecolor{ink}{HTML}{172A38}
\definecolor{accent}{HTML}{176866}
\definecolor{pale}{HTML}{EFF7F5}
\definecolor{muted}{HTML}{526570}
\hypersetup{colorlinks=true,urlcolor=accent,linkcolor=accent,
  pdftitle={How the proposed proof of Dirichlet GRH fits together},
  pdfauthor={Codex},pdfsubject={A mathematical reconstruction of the November 2026 announcement}}
\urlstyle{same}
\color{ink}
\setlength{\parindent}{0pt}
\setlength{\parskip}{5.5pt plus 0.5pt minus 0.5pt}
\setlength{\abovedisplayskip}{6pt plus 1pt minus 1pt}
\setlength{\belowdisplayskip}{6pt plus 1pt minus 1pt}
\setlength{\abovedisplayshortskip}{3pt}
\setlength{\belowdisplayshortskip}{4pt}
\setlist[enumerate]{leftmargin=1.55em,itemsep=3pt,topsep=2pt,parsep=0pt}
\pagestyle{fancy}
\fancyhf{}
\fancyhead[L]{\small\sffamily\color{muted}Dirichlet GRH}
\fancyhead[R]{\small\sffamily\color{muted}Mathematical reconstruction}
\fancyfoot[L]{\footnotesize\sffamily\color{muted}20 November 2026}
\fancyfoot[R]{\footnotesize\sffamily\color{muted}\thepage}
\renewcommand{\headrulewidth}{0.3pt}
\renewcommand{\headrule}{\hbox to\headwidth{\color{muted!45}\leaders\hrule height \headrulewidth\hfill}}
\newcommand{\heading}[1]{{\sffamily\bfseries\fontsize{22}{25}\selectfont #1\par}\vspace{3pt}}
\newcommand{\subhead}[1]{\par\addvspace{5pt}{\sffamily\bfseries\large\color{accent}#1\par}\nobreak\vspace{-2pt}}
\newcommand{\notes}[1]{{\small\sffamily\color{muted}#1}}
\newcommand{\repobase}{https://github.com/TomasOrtega/RH/blob/a6ac91a412dcbaab31801355ec5707024a71bf65/extracted/lean/src/}
\newcommand{\src}[3]{\href{\repobase#1.lean\#L#2}{#3}}
\newcommand{\aaDir}{Development/AutomorphicArithmetic/}
\newcommand{\chars}{Development/AutomorphicArithmetic/CharacterBoxes/}
\newcommand{\prim}{Development/AutomorphicArithmetic/PrimitiveCosets/}
\newcommand{\ded}{Development/AutomorphicArithmetic/DedekindLFunctions/}
\newcommand{\landau}{Development/AutomorphicArithmetic/DedekindLFunctions/LandauTheorems/}
\newcommand{\frames}{Formalism/SharpRestriction/RamificationFrobenius/TotientGaloisClasses/}
\newcommand{\note}[2]{\href{https://alpo.ge/remark.pdf\#page=#1}{#2}}
\newcommand{\R}{\mathbb R}
\newcommand{\C}{\mathbb C}
\newcommand{\F}{\mathbb F}
\newcommand{\Q}{\mathbb Q}
\newcommand{\Z}{\mathbb Z}
\newcommand{\re}{\operatorname{Re}}
\newcommand{\Res}{\operatorname*{Res}}
\newtcolorbox{keybox}{enhanced,breakable=false,sharp corners,boxrule=0pt,
  colback=pale,colframe=pale,borderline west={2pt}{0pt}{accent},
  left=9pt,right=9pt,top=5pt,bottom=5pt,before skip=5pt,after skip=4pt}
\begin{document}
\pdfbookmark[1]{The argument in one page}{overview}
{\sffamily\bfseries\fontsize{25}{28}\selectfont How the proposed proof of\\Dirichlet GRH fits together\par}
\vspace{2pt}
Dirichlet GRH says that every zero $\rho$ of a Dirichlet $L$-function in the strip $0<\re\rho<1$ has $\re\rho=\tfrac12$. A zero to the right would force a weighted sum over prime powers to be large somewhere in \emph{every} sufficiently late interval $[Y,Y^{3/2}]$. Levent Alpöge's construction is meant to produce arbitrarily late intervals on which the same sum stays small \emph{everywhere}. These statements give the contradiction. The difficult direction is producing the intervals with a bound arbitrarily close to the square-root scale. \notes{\note{3}{Announcement, pp.\ 3--4}.}

\subhead{The prime sum}
Write $\Lambda(p^a)=\log p$ for primes $p$ and integers $a\geq1$, and $\Lambda(n)=0$ otherwise. For $x\geq1$, define the weighted sum and its centered version by
\[
 R_\chi(x)=\sum_{n\leq x}\Lambda(n)\chi(n)(x/n-1),\qquad
 G_\chi(x)=
 \begin{cases}
 R_\chi(x)-x\log x+(1+\gamma)x,&\chi=\mathbf1\ \text{of conductor }1,\\
 R_\chi(x)+x\,L'(1,\chi)/L(1,\chi),&\chi\ \text{nonprincipal}.
 \end{cases}
\]
Here $\gamma$ is Euler's constant. For conductor one, $G_{\mathbf1}$ measures the fluctuation after the expected main term of the prime sum has been subtracted. The corrections remove the singularity at $s=1$ in the transform on page 2. Fix a primitive character $\psi$. Take the primitive characters inducing the powers of $\psi$, counting each distinct character once, and call this finite set $F_\psi$. For $h=\operatorname{ord}(\psi)$ the release uses
\[
 H_\psi(x)=G_{\mathbf1}(x)+\frac{1}{3h+1}
              \sum_{\chi\in F_\psi\setminus\{\mathbf1\}}\re G_\chi(x).
\]
For instance, the nontrivial character modulo $4$ has order two, so its family sum is $G_{\mathbf1}+\tfrac17\re G_{\chi}$. The family and its positive weights stay fixed as $Y$ varies. \notes{\src{\ded CharFamily}{10}{Prime sums} $\cdot$ \src{\frames ExplicitFormulaErrorOfHostConditions}{105}{cyclic family} $\cdot$ \src{\prim LevelFamily}{27}{weights}.}

\begin{keybox}
\textbf{\sffamily The required window estimate.}
For every primitive $\psi$, every $\tfrac12<\theta<1$, and every $Y_0\geq2$, there is a $Y\geq Y_0$ such that
\[
       H_\psi(x)\leq x^\theta\qquad\text{for every }Y\leq x\leq Y^{3/2}. \tag{W}
\]
\end{keybox}

\subhead{From the estimate to GRH}
Suppose $L(\rho,\psi)=0$ with $\tfrac12<\beta=\re\rho<1$. Choose $\tfrac12<\theta<\beta$. The window criterion on page 2 forces a violation of (W) in every sufficiently late window. This rules out $\rho$. A primitive zero with $\re\rho<\tfrac12$ in the open strip would give a zero of $L(1-\rho,\bar\psi)$ on the right by the functional equation. For an imprimitive character the extra Euler factors are nonzero when $\re s>0$, since $|\chi^*(p)p^{-s}|<1$ for its inducing primitive character $\chi^*$. The conductor-one case gives classical RH. \notes{\src{\chars ClaimStrip}{20}{Primitive reduction and reflection} $\cdot$ \src{Main}{4}{released conclusions}.}

\newpage
\pdfbookmark[1]{Why a zero forces a large prime sum}{zeros}
\heading{Why a zero forces a large prime sum}
\subhead{The smoothing gives each zero a pole}
Write $\widehat f(s)=\int_1^\infty f(x)x^{-s-1}\,dx$. For $\re s>1$, the contribution of a single ramp follows by integration,
\[
 \int_n^\infty (x/n-1)x^{-s-1}\,dx
   =n^{-s}\!\left(\frac1{s-1}-\frac1s\right)=\frac{n^{-s}}{s(s-1)} .
\]
Summing using the Euler product gives $\widehat R_\chi(s)=-(L'/L)(s,\chi)/[s(s-1)]$ by absolute convergence. For the conductor-one character this yields
\[
 \widehat G_{\mathbf1}(s)=
 -\frac{(\zeta'/\zeta)(s)}{s(s-1)}
 -\frac{1}{(s-1)^2}+\frac{1+\gamma}{s-1}.
\]
Indeed $-\zeta'/\zeta=1/(s-1)-\gamma+O(s-1)$ at $1$, so the two singular terms cancel. For nonprincipal $\chi$ the correction adds $(L'/L)(1,\chi)/(s-1)$ and again removes the pole at $1$. At a strip zero $\rho$ of multiplicity $r$, the continued transform has residue
\[
       \Res_{s=\rho}\widehat G_\chi(s)=-\frac{r}{\rho(\rho-1)}.
\]
Under Mellin inversion a pole at $\rho=\beta+it$ contributes a multiple of $x^\beta e^{it\log x}$. This explains why an exponent $\beta>\theta$ is detectable. Other zeros can interfere in the prime sum, so a single term does not yet locate a large value on the real $x$-axis.

Taking $\re G_\chi$ averages the transforms for $\chi$ and $\bar\chi$. At any common zero, all nonzero residues are positive multiples of $-1/[\rho(\rho-1)]$. Thus a zero of the target character survives the family average. At a real strip zero the residue is positive. \notes{\src{\landau MellinRamp}{10}{Ramp integral} $\cdot$ \src{\landau mellinGchiModOne}{115}{zeta centering} $\cdot$ \src{\landau poleData}{248}{family poles}.}

\subhead{Why control on windows suffices}
The release uses a Landau-type lemma. Set $H_F=\sum_{\chi\in F}a_\chi\re G_\chi$ for a fixed finite family with $a_\chi>0$ and the principal character represented at conductor one. For $0<\theta<1$ and $\lambda>1$, an upper bound $H_F(x)\leq Cx^\theta$ on arbitrarily late whole intervals $[Y,Y^\lambda]$ excludes every strip zero of every member with $\re\rho>\theta$.

Put $u=\log x$, $g(u)=H_F(e^u)$ and $\Phi(s)=\widehat H_F(s)$. A window becomes $[U,\lambda U]$. For $\sigma>2$, the normalized $k$th derivative $A_k(s)=(-1)^k\Phi^{(k)}(s)/k!$ is a Fourier integral at $s=\sigma+it$ with weight $w_k(u)=u^k e^{-\sigma u}/k!$. The triangle inequality gives
\[
 |A_k(\sigma+it)|+\re A_k(\sigma)
 \ \leq\ \int_0^\infty (|g(u)|+g(u))\,w_k(u)\,du.
\]
The integrand only sees the positive part of $g$. One chooses $\sigma$ and then $k\geq U$ so that the tails outside the window are small even using the crude bound $|g(u)|\ll e^{2u}$. For every $\eta>0$ this gives an upper bound $\ll e^{\eta U}(\sigma-\theta)^{-k}$ on the left along arbitrarily late windows.

A pole to the right of $\theta$ gives a conflicting lower bound. The source isolates a nearest pole at distance $d<\sigma-\theta$ from a suitable center $\sigma+it$, or treats a real pole using its positive residue. Its derivative contribution grows like $d^{-k-1}$ relative to an analytic remainder of strictly smaller exponential rate. Taking $\eta$ small gives the contradiction. The tail estimate and pole isolation are substantive parts of this analytic lemma. \notes{\src{\landau landauTails}{80}{Tails} $\cdot$ \src{\landau landauWindow}{83}{pole argument} $\cdot$ \src{\landau ZeroFreeOfRpowWindow}{11}{window criterion}.}

\newpage
\pdfbookmark[1]{What the theta lifts contribute}{lifts}
\heading{How the theta lifts enter the estimate}
\subhead{An exact calculation in the finite-field model}
Let $e_q$ be a nontrivial additive character of $\F_q$, let $t:\F_q\to\C$, and put $M=\sum_x t(x)$. On $k\geq2$ variables take the path quadratic form $Q(\mathbf x)=\sum_{j=1}^{k-1}x_jx_{j+1}$ and $T(\mathbf x)=\prod_j t(x_j)$. Summing over the additive frequency first, or separating its zero term first, gives
\[
 q\!\sum_{Q(\mathbf x)=0}T(\mathbf x)
 =\sum_{a\in\F_q}\sum_{\mathbf x} T(\mathbf x)e_q(aQ(\mathbf x))
 =M^k+\sum_{a\ne0}\sum_{\mathbf x}T(\mathbf x)e_q(aQ(\mathbf x)).
\]
This recovers the announcement's $qN=M^k+R$. There is a concrete saving for each $a\ne0$. The matrix $(q^{-1/2}e_q(axy))_{x,y}$ is unitary by character orthogonality. Each of the $k-1$ edges contributes a factor $\sqrt q$. Multiplication by $t$ has operator norm at most one when $|t|\leq1$, and the two endpoint vectors each have norm at most $\sqrt q$. Hence each inner sum has absolute value at most $q^{(k+1)/2}$. Controlling $M$ also requires information about the constrained sum on the left. The announcement attributes a sign there to geometry. \notes{\note{3}{Announcement, footnote 3}.}

\subhead{From vanishing order to arithmetic divisibility}
In Stepanov's method a nonzero auxiliary function vanishing to order $m$ at each of $N$ points has at least $mN$ zeros counted with multiplicity. A degree bound then controls $N$. The proposed arithmetic analogue builds an auxiliary determinant from periods of automorphic forms. A period is an integral of a form over a subgroup. The determinant and the branching of representations are supposed to enforce simultaneous vanishing or divisibility conditions.

One local prototype in the note is the congruence $D^p\equiv0\bmod p$ for an invariant derivation on a CM elliptic curve at a good inert prime. It gives a source of divisibility in high theta derivatives. The note reports multiplicity two on $\mathrm O(2,n)$, then iterates theta lifts between definite unitary groups $U(V_0),\ldots,U(V_m)$. A theta lift integrates a form against a theta kernel to obtain a form on the next group. On the chosen input lattice, each lift carries forward the valuation already earned and adds a local gain less its denominator cost. In the comparison with primes, the same signal occurs once for each of the $m$ lifts. For a prime sum of length $x$ the release uses
\[
                x^{1/2-\varepsilon}\ \leq\ m\ \leq\ x^{1/2+\varepsilon/2}.
\]
After division by the repeated signal, the leading costs are of order $x/m$ and $m$, up to powers of $x^\varepsilon$. This explains the choice of scale. \notes{\note{3}{Announcement, items 2 and 7--14} $\cdot$ \src{Formalism/SizeWindowBrackets/SChain}{53}{lower scale} $\cdot$ \src{\landau ChCore}{343}{upper scale}.}

\subhead{How the cost of a lift enters}
For a nonzero algebraic integer $a$ in a number field $K$, finite divisibility and archimedean size are linked by the norm identity
\[
 \sum_{\mathfrak p}v_{\mathfrak p}(a)\log N\mathfrak p
 =\sum_{\sigma:K\hookrightarrow\C}\log|\sigma(a)|.
\]
A slope comparison applies this principle to lattices. For fixed auxiliary member $u$, the periods at different Hecke twists factor through one bilinear form. At each prime ideal $\mathfrak p$, let $q_{\mathfrak p}$ be the certified integer lower bound on the valuations of its entries on the chosen local lattices. Multiplying the input lattice by $\prod_{\mathfrak p}\mathfrak p^{-q_{\mathfrak p}}$ still gives integral pairings. This is the common division in the announcement. The archimedean norm and the remaining congruence losses must also be bounded. \notes{\note{4}{Announcement, items 12--18} $\cdot$ \href{https://github.com/robbiemu/dirichlet-grh-sketch/blob/main/lift-estimate.pdf}{worked lift estimate, pp.\ 1--3}.}

\newpage
\pdfbookmark[1]{The arithmetic estimate and its status}{arithmetic}
\heading{The arithmetic estimate behind the claim}
\subhead{Two inputs must hold at the same scales}
Fix $\psi$ and its cyclic character family. Write $S_\varepsilon$ for the scales of a chosen auxiliary family where the field, character data and chain estimates can be used together.
\begin{enumerate}
\item \textbf{Construct the families.} There exist constants $e_D,e_L$ and $\varepsilon_0>0$ such that, for each $0<\varepsilon<\varepsilon_0$, one can choose a compatible family. Its admissible scales contain an entire interval $[Y,Y^{3/2}]$ above each prescribed starting point. On all sufficiently large scales in this set the divisibility, level and norm estimates apply to the same chain.
\item \textbf{Compare with the prime sum.} With these data, at every sufficiently large $x\in S_\varepsilon$ the arithmetic comparison yields
\[
                      H_\psi(x)\ \leq\ x^{1/2+(e_D+e_L+7)\varepsilon}. \tag{1}
\]
The connection with primes is finite Fourier inversion. In a finite abelian group $A$,
\[
  \mathbf1_{g=a}=\frac1{|A|}\sum_{\xi\in\widehat A}\overline{\xi(a)}\,\xi(g).
\]
At an unramified prime, inserting $g=\mathrm{Frob}_p^r$ converts a condition on a Frobenius class into character weights on $p^r$. The arithmetic side bounds the excess of such weighted counts over their main term. Averaging the resulting identities gives (1). \notes{\src{\chars ZeroFreeAllOfTopG}{12}{Compatible families} $\cdot$ \src{\landau ChBalance}{42}{orthogonality} $\cdot$ \src{\landau ChTop}{37}{excess bound} $\cdot$ \src{\prim CoreGm}{12}{comparison}.}
\end{enumerate}
The release takes $e_D=2$ and $e_L=0$ before $\varepsilon$ is chosen. Thus (1) reads $H_\psi(x)\le x^{1/2+9\varepsilon}$. Choose $0<\varepsilon<\min\{\varepsilon_0,(\theta-1/2)/9\}$ and a whole window beyond the comparison threshold to obtain (W). The four remaining error terms and their exponents are calculated separately in the companion. \notes{\src{Development/AutomorphicArithmetic/ThetaCorrespondence/ThetaPairLifts/ThirdPlaceSeries/TopShape}{52}{Choice of constants} $\cdot$ \src{\frames TopWin}{68}{exponent choice} $\cdot$ \href{https://github.com/robbiemu/dirichlet-grh-sketch/blob/main/lift-estimate.pdf}{worked lift estimate, p.\ 4}.}

\subhead{An example of what the construction has to prove}
The announcement mentions a weakening of Vandiver. Its role can be stated quantitatively. For specified auxiliary cyclotomic subfields $E$, write $b_{\ell,j}(E)$ for the $\Z_\ell$-length of the locally trivial subgroup of the relevant $H^1$ with coefficients $\Q_\ell/\Z_\ell(-j)$. On certain auxiliary windows the release requires these lengths to be finite and, for every finite set $P$ of primes, requires
\[
 \sum_{\ell\in P}(\log\ell)\!\!\sum_{1\leq j<\lfloor x^{401}\rfloor}
             b_{\ell,j}(E)\ \leq\ \exp(x^\varepsilon/2).
\]
The source specifies the fields and the allowed levels in \texttt{VWinAt}, with its own auxiliary intervals $[Z,Z^2]$. This bound controls the contribution of the surviving cohomology classes to the cost of division. The announcement also singles out a Behrend construction in choosing its Iwasawa tower. \notes{\src{Formalism/GaloisCohomology/Statement}{67}{Cohomology group} $\cdot$ \src{\chars ExistsPageN}{12}{field and window conditions} $\cdot$ \src{\chars DivinNOfHypV}{17}{use in divisibility}.}

\subhead{What was reproduced}
We worked out the finite Fourier identity, the Mellin residues, and the implication from (W) to Dirichlet GRH using the analytic window criterion. For the lifts we reconstructed the local valuation induction, the common lattice division, the main-term cancellations and the final error calculation. Independent verification of the automorphic period and norm inputs, the uniform construction on the same windows, and the full Lean build remains open here. \notes{\src{\chars MainConclusion}{7}{Released assembly} $\cdot$ \src{Statements}{3}{formal statements}.}
\end{document}
